% CLAUDE REPORT::B3-CONSERVATION-0001::2026-11-12
\documentclass[11pt]{article}
\usepackage[margin=1in]{geometry}
\usepackage{amsmath,amssymb,amsthm}
\usepackage{booktabs}
\usepackage[hidelinks]{hyperref}
\newtheorem{theorem}{Theorem}
\newtheorem{lemma}[theorem]{Lemma}
\newtheorem{corollary}[theorem]{Corollary}
\theoremstyle{definition}
\newtheorem{definition}[theorem]{Definition}
\newcommand{\Wbar}{\overline W}
\newcommand{\SAT}{\mathrm{SAT}}
\newcommand{\UNSAT}{\mathrm{UNSAT}}
\newcommand{\poly}{\operatorname{poly}}
\title{The B3 bridge pays for itself:\\
a conservation theorem for the worst-case minimum-derivation programme}
\author{Claude scheduled research run (worst-case derivation directive)}
\date{12 November 2026}
\begin{document}
\maketitle

\section*{Status}
\textbf{P versus NP remains OPEN.} Theorems~1, 3, 4 and Corollary~2 below are
\textbf{PROVED/EXACT} from the stated definitions and cited known theorems.
Theorem~5 is a \textbf{KNOWN THEOREM} pattern (Cook 1975; Krajíček--Pudlák 1989)
restated in the project's notation. The bridge B3 for any fixed system with an
unproved lower bound, and every lower bound for Frege/Extended Frege, remain
\textbf{OPEN}. Nothing here is a separation.

\section{Setting (from the 2026-11-02 shared directive)}
Fix a sound, complete, polynomial-time-checkable refutation system $R$ with
nonnegative node weights. For UNSAT $F$,
\[
W_R(F)=\min_{\pi:F\vdash_R\bot}\operatorname{cost}_R(\pi),\qquad
\Wbar_R(N)=\max_{F\in\UNSAT,\ |\langle F\rangle|\le N} W_R(F).
\]
The inner minimum grants free optimal choice of the derivation (the user's
point that search can only add work is accepted: this is proof \emph{length},
not proof \emph{search}). The primary bridge is
\[
\textbf{B3}_R:\qquad \SAT\in \mathrm P\ \Longrightarrow\ \exists C,k\ \forall N:\ \Wbar_R(N)\le C(N+1)^k .
\]
Write $\mathrm{LB}_R$ for ``$\Wbar_R$ is superpolynomial'' and
$\mathrm{PB}_R$ for ``$\Wbar_R$ is polynomially bounded'' ($\mathrm{PB}_R=\neg\mathrm{LB}_R$).

\section{Barrier-transfer preflight}
\textbf{Obstacle 1:} B3 for a fixed $R$ (which algorithm-to-proof translation?).
\textbf{Obstacle 2:} a superpolynomial $\Wbar_R$ for an $R$ strong enough to
make B3 plausible (shared DAG reasoning, extension/auxiliary variables).
Literature checked this run: Cook--Reckhow (JSL 1979); Cook (STOC 1975, PV and
EF); Krajíček--Pudlák (JSL 1989, optimal proof systems and reflection);
Haken (TCS 1985); Urquhart (JACM 1987); Pudlák (JSL 1997, cutting planes);
Ajtai (1988), Krajíček--Pudlák--Woods and Pitassi--Beame--Impagliazzo (1995),
Håstad (depth up to $\sim\log n/\log\log n$); Razborov (1998, PC);
de~Rezende--Engström--Ghannane--Risse, ECCC TR26-078 / arXiv:2604.28172
(2026: superpolynomial lower bounds for \emph{tree-like} semantic Frege with
bounded line size); Davis--Robere, ECCC TR26-055 (2026).
Transfer domain: logic (reflection principles / Gödel). Classification:
\emph{directly transferable theorem} for Theorems~3 and~5; the 2026 tree-like
result is a \emph{useful analogy requiring a new theorem} for the project's
DAG-like hypergraph calculus (DAG sharing is exactly what it does not cover).

\section{Results}

\begin{theorem}[Truth-table of B3; PROVED/EXACT]
For every fixed $R$: \ $\mathrm{B3}_R \iff (\SAT\notin\mathrm P)\ \vee\ \mathrm{PB}_R$.
In particular $\mathrm{P}\ne\mathrm{NP}\Rightarrow \mathrm{B3}_R$ for every $R$, and
$\mathrm{B3}_R\wedge \mathrm{LB}_R\Rightarrow \mathrm P\neq\mathrm{NP}$.
\end{theorem}
\begin{proof}
B3 is a material implication whose consequent is $\mathrm{PB}_R$; the
consequent's constants do not depend on any SAT algorithm. Expand
$A\Rightarrow B\equiv\neg A\vee B$.
\end{proof}

\begin{corollary}[Conservation for systems with known lower bounds; PROVED/EXACT]
If $\mathrm{LB}_R$ is a theorem, then $\mathrm{B3}_R\iff \mathrm P\neq\mathrm{NP}$.
This applies (KNOWN THEOREM inputs) to Resolution, Cutting Planes,
bounded-depth Frege of any fixed depth, Polynomial Calculus, and every
restriction of them, including any multisign hypergraph calculus whose
derivations are p-simulated by one of these.
\end{corollary}
\begin{proof}
With $\mathrm{PB}_R$ false, Theorem~1 leaves $\mathrm{B3}_R\iff\SAT\notin\mathrm P$.
The listed systems have unconditional exponential lower bounds
(PHP / Tseitin / clique--colouring families), so $\mathrm{LB}_R$ holds;
lower bounds transfer down along p-simulation.
\end{proof}

\noindent\emph{Consequence for the programme.} For every currently analysable
$R$, ``proving the weakest bridge'' is \emph{literally the full P$\ne$NP
problem}; B3 is not weaker than the goal there. The parity-recognizer
observation (B1 false for Resolution) was the pointwise shadow of this; the
corollary shows that even the worst-case B3 cannot be a stepping stone for
these $R$. Any proof of B3 must therefore be \emph{uniform/constructive}: it
must turn an arbitrary hypothetical polynomial SAT algorithm into short
$R$-refutations without first deciding P vs NP.

\begin{theorem}[Optimal systems make B3 free; PROVED/EXACT from Cook--Reckhow]
If $R$ is optimal (for every proof system $Q$ there is a polynomial $q$ with
$\Wbar_R(N)\le q(N+\Wbar_Q(N))$, unit-like weights polynomially related to
proof bit-length), then $\mathrm{B3}_R$ holds unconditionally, and moreover
$\mathrm{LB}_R\iff \mathrm{NP}\neq\mathrm{coNP}$.
\end{theorem}
\begin{proof}
If $\SAT\in\mathrm P$ via $A$ then $Q_A$ (proof of $F$ is the empty string,
verifier runs $A$) is a proof system with $\Wbar_{Q_A}(N)=O(1)$ in proof length,
so $\Wbar_R$ is polynomial. For the equivalence: NP$=$coNP iff some
p-bounded proof system exists (Cook--Reckhow); optimality transfers boundedness
to $R$, and conversely $\mathrm{PB}_R$ makes $R$ itself p-bounded.
\end{proof}

\begin{theorem}[No free middle; PROVED/EXACT]
Let $\mathcal S$ be the set of systems $R$ for which $\mathrm{B3}_R$ is provable
in ZFC \emph{without} proving $\mathrm{P}\ne\mathrm{NP}$ or $\mathrm{PB}_R$,
and $\mathcal L$ those for which $\mathrm{LB}_R$ is provable. Any $R\in\mathcal
S\cap\mathcal L$ yields a ZFC proof of P$\neq$NP. Hence exhibiting such an $R$
is equivalent in difficulty to the original problem, and currently
$\mathcal L$ contains only systems for which (Corollary~2) $\mathrm{B3}_R$
is equivalent to P$\neq$NP.
\end{theorem}
\begin{proof}Immediate from Theorem~1.\end{proof}

\begin{theorem}[Reflection route to B3; KNOWN THEOREM pattern]
Let $R$ p-simulate Extended Frege (EF). Suppose $\SAT\in\mathrm P$ is
witnessed by a clocked algorithm $A$ such that the propositional translations
$\|\mathrm{Ref}_A\|_n$ (``if the computation of $A$ on input $\langle F\rangle$
of length $n$ outputs UNSAT then $F(x)$ is false'') have $R$-proofs of size
$\poly(n)$. Then $\Wbar_R(N)\le\poly(N)$. For $R=$EF this hypothesis holds
whenever PV (Cook's theory) proves that $A$ is correct.
\end{theorem}
\begin{proof}[Proof sketch]
Given UNSAT $F$: (i) instantiate the short proof of $\|\mathrm{Ref}_A\|_n$ at
the fixed input bits $\langle F\rangle$; (ii) EF evaluates the poly-size
computation circuit of $A$ on constant inputs in poly size, deriving the
``outputs UNSAT'' antecedent; (iii) modus ponens gives $\neg F(x)$; (iv) with
the clauses of $F$ derive $\bot$. Each step is polynomial. The PV clause is
Cook's translation theorem (PV-provable $\forall\Pi^b_1$ statements have
poly-size EF proofs).
\end{proof}

\noindent\textbf{Exact remaining gap for strong $R$ (OPEN).} B3$_{\mathrm{EF}}$
follows from: \emph{if P$=$NP then some polynomial SAT algorithm has
PV-provable (or EF-short) correctness}. A world where P$=$NP holds but no
correct algorithm's correctness has short EF proofs is not excluded by any
known theorem; this is exactly the gap. Simultaneously $\mathrm{LB}_{\mathrm{EF}}$
(indeed even for Frege) is a famous open problem; the 2026 tree-like semantic
bounded-line result does not cover DAG-like proofs or extension variables.

\section{Hostile controls and red-team audit}
\begin{itemize}
\item \emph{Does Theorem~1 ``refute'' B3?} No. It is a CONDITIONAL-BRIDGE
  classification: B3 is true in every P$\ne$NP world; proving it for
  Resolution-like $R$ simply \emph{is} proving P$\neq$NP.
\item \emph{Free best sequence.} Accepted throughout; all results concern
  proof length. Granting free search does not affect Theorem~1.
\item \emph{Relativization.} Theorem~1 is purely logical; Theorem~3 relativizes;
  Theorem~5 is non-relativizing only through PV provability.
\item \emph{Encoding size vs.\ variables.} All bounds are in encoded length $N$.
\item \emph{Pathological systems} (LSPV-1/395 PACK/HALT): optimality in
  Theorem~3 is about proof bit-length; node-weight systems hiding evidence
  are excluded by the bit-length-related weight convention stated there.
\end{itemize}

\section{Corpus reviewed / duplication audit}
Read: \texttt{AGENTS.md}; the 2026-11-02 worst-case refocus handoff;
LSPV-1/335 (MU core, generator barrier) and LSPV-1/395 (noncomputable
worst-case node profile); Codex 2026-11-11 SAT polynomial assumption audit
(header and model); earlier Claude world-ledger/T-PORT/one-tape notes;
GROKPROOF log list. Repository search for ``vacuous''/reflection/Krajíček--Pudlák
in the B3 context found no prior statement of Theorem~1/Corollary~2 as a
conservation law for the bridge, nor of Theorem~3 for B3; these are elementary
and \emph{not claimed as mathematically novel}, only newly recorded. LSPV-1/335
lemma~2 (generator barrier) is consistent: it shows a \emph{universal pointwise}
transfer fails for explicit families; Theorem~1 concerns the worst-case
conditional bridge.

\textbf{What changed relative to root main:} no edit to root \texttt{main.tex}
(this run is not an LSPV ownership role). \textbf{Adopted:} directive
definitions, MU-core lemma. \textbf{Rejected:} treating B3 as strictly weaker
than P$\ne$NP for Resolution-type $R$. \textbf{Stale/overstated claims found:}
any description of B3 as an ``easier first step'' for systems with known lower
bounds is overstated by Corollary~2.

\section{Recommended next targets}
\begin{enumerate}
\item (Lead/Builder) Fix $R=$ EF$+$hypergraph rules; pursue the PV-provability
  form of B3 as a precise CONDITIONAL, and look for a world (oracle or
  model of PV) separating ``P$=$NP'' from ``P$=$NP with PV-provable witness''.
\item (Certificates) Seek DAG-like lower bounds for the project's hypergraph
  calculus beyond tree-like semantic bounded-line systems (arXiv:2604.28172 as
  the closest current technique).
\item (Red team) Check every proposed $R$ against Corollary~2 first: if a lower
  bound for $R$ is already known, B3$_R$ is the full problem.
\end{enumerate}

\begin{thebibliography}{9}
\bibitem{CR} S.~Cook, R.~Reckhow, The relative efficiency of propositional proof systems, JSL 44 (1979).
\bibitem{C75} S.~Cook, Feasibly constructive proofs and the propositional calculus, STOC 1975.
\bibitem{KP} J.~Krajíček, P.~Pudlák, Propositional proof systems, the consistency of first order theories and the complexity of computations, JSL 54 (1989).
\bibitem{H} A.~Haken, The intractability of resolution, TCS 39 (1985).
\bibitem{P97} P.~Pudlák, Lower bounds for resolution and cutting plane proofs and monotone computations, JSL 62 (1997).
\bibitem{R98} A.~Razborov, Lower bounds for the polynomial calculus, Comput.\ Complexity 7 (1998).
\bibitem{dR26} S.~de~Rezende, A.~Engström, Y.~Ghannane, K.~Risse, Superpolynomial length lower bounds for tree-like semantic proof systems with bounded line size, ECCC TR26-078, arXiv:2604.28172 (2026).
\end{thebibliography}
\end{document}
% END_OF_VERIFIED_REPORT::CLAUDE::B3-CONSERVATION-0001
